% ─────────────────────────────────────────────────────────────────────────────
% Supplementary：從正文 4.5 移出（使用者 2026-10-26 決定正文不放）。表格由 build_paper_results.py 產生（tab_analysis.tex）。
% 放進 supplementary 時，把 \subsection 改成 supplementary 的章節層級。
% ─────────────────────────────────────────────────────────────────────────────
\subsection{Where \method helps and where it fails}
\label{sec:analysis}

\providecommand{\pend}{\textcolor{red}{?}}

\begin{table}[t]
  \centering
  \footnotesize
  \setlength{\tabcolsep}{2pt}
  \begin{tabular}{l cc cc cc cc}
    \toprule
    & \multicolumn{2}{c}{HR-4K} & \multicolumn{2}{c}{HR-8K} & \multicolumn{4}{c}{V$^*$ Bench} \\
    \cmidrule(lr){2-3}\cmidrule(lr){4-5}\cmidrule(lr){6-9}
    Method & S & C & S & C & Attr. & Pos. & Hit & Miss \\
    \midrule
    Uncompressed & 77.0 & 47.0 & 66.0 & 47.0 & 85.2 & 81.6 & 86.8 & 76.4 \\
    \midrule
    FastV & 73.0 & \textbf{45.0} & 63.0 & \textbf{45.0} & 80.0 & 77.6 & 83.1 & \textbf{69.1} \\
    SparseVLM & 75.0 & 44.0 & 66.0 & 41.0 & 80.0 & \textbf{78.9} & 85.3 & 65.5 \\
    VisionZip & 62.0 & 36.0 & 59.0 & 30.0 & 75.7 & 61.8 & 73.5 & 61.8 \\
    IF-Prune & \textbf{76.0} & 39.0 & \textbf{68.0} & 37.0 & \textbf{85.2} & 76.3 & 86.8 & \textbf{69.1} \\
    \rowcolor{oursbg}
    \method (ours) & 71.0 & 42.0 & 61.0 & 43.0 & 77.4 & 77.6 & \textbf{88.2} & 50.9 \\
    \bottomrule
  \end{tabular}
  \caption{\textbf{Where \method helps and where it fails} ($B=1536$). HR-Bench is split into single- (S) and cross-instance (C) questions (CircularEval), V$^*$ Bench into attribute and relative-position questions. Hit / Miss split V$^*$ Bench by whether the detail tiles selected by \method cover the target (136 hit, 55 miss); every method is evaluated on the same split.}
  \label{tab:analysis}
\end{table}


\cref{tab:analysis} breaks down the high-resolution benchmarks by question type. On questions that
involve several objects or the global arrangement of the scene (HR-Bench cross-instance, V$^*$ Bench
relative position), \method is on par with or better than the baselines, \eg 43.0 \vs 37.0 for IF-Prune
on HR-Bench 8K cross-instance. Its gap is concentrated on questions about a single small object
(HR-Bench single-instance, V$^*$ Bench attributes).

V$^*$ Bench provides target boxes, which tell us why. When the selected detail tiles cover the target
(136 of 191 questions), \method answers 88.2\% correctly, more than any baseline and more than the
uncompressed model (86.8\%), since the target is seen at native resolution without the distraction of
the other tiles. When the target is missed, accuracy drops to 50.9\%. The misses come from the prior,
not from boundary scoring: most missed targets are ranked outside the scoring window, while targets
that enter the window are selected 93--95\% of the time ($B=1024$ and $2048$). The prior is computed
on the coverage tiles, in which more than half of the V$^*$ targets are smaller than one ViT patch
(14\,px), so a small budget can leave them out of the window. Methods
that encode every tile do not have this failure mode, but pay $\mathcal{O}(N)$ for it.

